\clearpage\thispagestyle{empty}
\pdfbookmark[0]{Sources for the Outline of Orbit Determination (1809)}{summarische1809-sources}
\begin{center}\textbf{Sources for the Outline of Orbit Determination (1809)}\end{center}
\noindent This AI translation of Carl Friedrich Gauss's German paper follows \href{https://gdz.sub.uni-goettingen.de/id/PPN236007467}{\textit{Werke}, vol.~VI} (Göttingen, 1874), pp.~148--165. It includes all nine articles, Bernhard von Lindenau's signed introductory note and Gauss's four notes. The \href{https://www.digitale-sammlungen.de/en/view/bsb10538610?page=203}{first printing} appeared as article XVII in the September 1809 issue of Franz Xaver von Zach's \textit{Monatliche Correspondenz zur Beförderung der Erd- und Himmels-Kunde}, vol.~20, pp.~197--224.
\par\smallskip\noindent Gauss's reference to Wilhelm Olbers's \textit{Abhandlung über die leichteste und bequemste Methode die Bahn eines Cometen aus einigen Beobachtungen zu berechnen} (Weimar, 1797) is to \href{https://www.e-rara.ch/zut/content/pageview/449839}{p.~45, §38}. For the integration formulas, see the postscript to Roger Cotes's \textit{De methodo differentiali Newtoniana}, in \textit{Harmonia mensurarum}, edited and expanded by Robert Smith (Cambridge, 1722), \href{https://www.e-rara.ch/zut/content/pageview/1256011}{pp.~32--33}. The first printing refers to \href{https://www.e-rara.ch/zut/content/pageview/142894}{p.~87} of Gauss's \textit{Theoria motus corporum coelestium in sectionibus conicis solem ambientium} (Hamburg, 1809), where Article~83 begins. The collected text cites that article.
\par\smallskip\noindent A separate \href{https://www2.karlin.mff.cuni.cz/~tuma/Aplikace21/ALL_Gauss_1809_Summarische.pdf}{English version}, with no translator named, is available from Jiří Tůma's Aplikace21 course at Charles University.
\par\smallskip\noindent Translation and typesetting produced by AI.
\clearpage
